\scanpage{211}
% L1--4 (haut de la page)
that the automorphisms of such a $G_\nu$ ($\simeq A_{g,\nu}$) which fix the class $U_\nu$ i.e. which \unsure{transform|transport} $H_\nu$ into a conjugate, \unsure{are} inner, and that the center of $H_\nu$ is $\{1\}$.

% L5--9
In fact, $U_\nu$ can also be interpreted as a homogeneous space under $G_\nu^\circ$, and not only under $G_\nu$, more generally we have that $G_\nu^\circ$ is transitive on $U_\nu$, and \unsure{even} on $\mathcal P_{\nu'}(U_\nu)$ for every $\nu'\in\mathbb N^*$.

% L10--14, formule (30)
Returning to the model situation with $U_{g,\nu}$, we find
\[
\tag{30}
\begin{aligned}
U_{g,\nu}
&\simeq A_{g,\nu}^{\unsure{\deleted{\circ}|\bullet}}/
\underbrace{A_{g,(\nu,\nu+1)}}_{\substack{\text{\unsure{$\parallel$ definition}}\\ A_{g,\nu}\cap A_{g,\nu+1}}}\\[1.2em]
&\simeq A_{g,\nu}^{\circ}/
\underbrace{B_{g,(\nu,\nu+1)}}_{\substack{\text{\unsure{$\parallel$ (definition)}}\\
 A_{g,\deleted{\unsure{(\nu,}}\unsure{\{\nu+1\}}}\cap A_{g,\nu}^{\circ}
 =\{u\in A_{g,\nu}^{\circ}\mid u.a_{g,\nu}=\unsure{a_{g,\nu}}\}}}.
\end{aligned}
\]
% Exposant du premier numérateur : petite tache d'encre pleine (cercle noirci ou point), ni « ! » ni « * ».

% L15--18, bas de la page
We thus have \deleted{\unsure{some}} \deleted{\unsure{inclusions}} \inserted{by inverse} diagram of inclusions of subgroups of $A_{g,\nu}$ :
% « par inv. » est écrit en petit sous le mot barré « inclusions ».
\deleted{\diagram{\begin{tikzcd}[ampersand replacement=\&,column sep=1.7em,row sep=2.7em]
\unsure{A_{g,\illegible{1}}} \arrow[r,phantom,"\unsure{\supset}" description] \&
\unsure{A_{g,\nu}} \arrow[r,phantom,"\supset" description] \&
A_{g,\nu}^{!} \arrow[r,phantom,"\unsure{\supset}" description] \&
A_{g,\nu}^{\circ} \arrow[r,phantom,"\supset" description] \&
B_{g,\nu} \arrow[r,phantom,"\supset" description] \arrow[dl,hook,bend left=15,line width=0.8pt] \&
B_{g,\nu}^{\circ} \\
{} \& {} \& {} \& A_{g,(\nu,\nu+1)} \arrow[ull,hook,no head,bend left=10] \& {} \& {}
\end{tikzcd}}}
\margin{210}

\scanpage{212}
\mspage{106}
% L1--5, diagramme (31)
\diagram{\begin{tikzcd}[ampersand replacement=\&,column sep=2.1em,row sep=3.2em]
|[yshift=-1.4em]| {(31)} \&
A_{g,\nu} \arrow[rr,bend left=24,no head,"\Gamma_{g,\nu}"] \&
A_{g,\nu}^{!} \arrow[l,hook,"\mathrm{inv.}"',"\unsure{\mathfrak S_{\nu}}"] \&
A_{g,\nu}^{\circ} \arrow[l,hook,"\mathrm{inv.}"',"\Gamma_{g,\nu}^{!}"] \& {} \\
\left[A_{g,\nu+1}\,\hookleftarrow\right] \&
A_{g,(\nu,\nu+1)} \arrow[u,hook] \arrow[rr,bend right=45,no head,"\Gamma_{g,\nu}"'] \&
A_{g,\nu+1}^{!} \arrow[u,hook] \arrow[l,hook,"\substack{\mathrm{inv.}\\\unsure{\mathfrak S_{\nu}}}"] \&
B_{g,\nu} \arrow[u,hook] \arrow[l,hook,"\substack{\mathrm{inv.}\\\Gamma_{g,\nu}^{!}}"] \arrow[d,equal] \&
B_{g,\nu}^{\circ}=(A_{g,\nu+1}^{!})^{\circ}=A_{g,\nu+1}^{\circ} \arrow[l,hook,"\mathrm{inv.}"] \\[-1.2em]
{} \& {} \& {} \& \left[(A_{g,\nu}^{\circ})_{\unsure{a_{g,\nu}}}\right] \& {}
\end{tikzcd}}

% Marge gauche, à hauteur de L4--13
\margin{\uline{N.B.} $A_{g,(\nu,\nu+1)}$, \unsure{stabilizer} \illegible{2} in $(A_{g,\nu}=\operatorname{Aut}(U_{g,\nu}))$ of the \unsure{marked point} $a_{g,\nu}$ \unsure{of} \illegible{1}, \unsure{deserves} a simpler notation, for example $A_{g,\nu}^{\bullet}$ (the exponent $\bullet$ \unsure{here} to \unsure{emphasize} \illegible{3} \uline{\unsure{pointed|pointed}} in $U_{g,\nu}$).}

% L6--9
The two squares are cartesian, and the horizontal inclusions \inserted{(except the one in \unsure{brackets} in the $2^{\text{nd}}$ row)} are inclusions of invariant subgroups \inserted{invariant even} (in the largest group, $A_{g,\nu}$ resp. $A_{g,(\nu,\nu+1)}$).

% L10--15
[The equality $B_{g,\nu}^{\circ}=A_{g,\nu+1}^{!\,\circ}$ \deleted{\unsure{by}} (which refines the trivial inclusion $B_{g,\nu}^{\circ}\subset A_{g,\nu+1}^{!\,\circ}$ by the reverse inclusion $A_{g,\nu+1}^{!\,\circ}\subset B_{g,\nu}^{\circ}$, equivalent to $A_{g,\nu+1}^{!\,\circ}\subset B_{g,\nu}$) comes from the inclusion
\[
A_{g,\nu+1}^{!\,\circ}\inserted{\unsure{=}A_{g,\nu+1}^{\circ}}\subset A_{g,\nu}^{!\,\circ}=A_{g,\nu}^{\circ}
\quad\text{(whence : }A_{g,\nu+1}^{\circ}\subset A_{g,\nu+1}\cap A_{g,\nu}^{\circ}=B_{g,\nu}\text{).}
\]
% L'ajout interlinéaire est écrit au-dessus du premier \subset, relié à lui par un trait vertical que coupent deux petits traits.
]

% L16--19
The three vertical inclusions define three homogeneous spaces, and \deleted{$\unsure{A_g}$} inclusion homomorphisms between them which are \uline{bijective}
\[
(\unsure{U_{g,\nu}}\simeq)\quad
A_{g,\nu}/A_{g,(\nu,\nu+1)}
\xleftarrow{\sim} A_{g,\nu}^{!}/A_{g,\nu+1}^{!}
\xleftarrow{\sim} A_{g,\nu}^{\circ}/B_{g,\nu},
\]
% Le U de (U_{g,\nu}\simeq) est repassé à l'encre forte, non barré.
\inserted{\unsure{with} (transitivity !)}
and \deleted{\unsure{likewise by}} \deleted{\unsure{arrows} \illegible{2}} \inserted{the quotient groups \unsure{denoted} $\mathfrak S_\nu$ and $\Gamma_{g,\nu}^{!}$} defined by the horizontal inclusions of the first row, are isomorphic via the vertical inclusions to the corresponding quotients in the second row. Thus the extension of groups
\[
(12)\qquad 1\longrightarrow\Gamma_{g,\nu}^{!}\longrightarrow\Gamma_{g,\nu}\longrightarrow\unsure{\mathfrak S_\nu}\longrightarrow1
\]
% L25--26, fin de page
can \deleted{also} be deduced equally well from
\margin{211}

\scanpage{213}
% L1--5, formule (32)
the first row of (31), or from the first three terms of the second, in particular
\[
\tag{32}
\left\{\begin{aligned}
\unsure{\mathfrak S_\nu}&\simeq A_{g,(\nu,\nu+1)}/A_{g,\nu+1}^{!}\\
\Gamma_{g,\nu}^{\unsure{\bullet|\,}}&\simeq A_{g,(\nu,\nu+1)}/B_{g,\nu}\\
\Gamma_{g,\nu}^{!}&\simeq A_{g,\nu+1}^{!}/B_{g,\nu}.
\end{aligned}\right.
\]
% Au-dessus de \Gamma_{g,\nu} (deuxième ligne) : point d'encre plein, sans cercle ; peut-être une tache.

% L6--12
Thus the \inserted{\unsure{canonical}} torsor with group $A_{g,(\nu,\nu+1)}$ \inserted{$=(A_{g,\nu})_{\unsure{a_{g,\nu}}}$} over \inserted{the homogeneous space} $U_{g,\nu}$, and \inserted{the one} of the group $A_{g,\nu+1}^{!}$ ($\subset A_{g,(\nu,\nu+1)}$) which gives rise to it, is in either case deduced, by extension of the structure groups, from a torsor with group $B_{g,\nu}$ (whose fiber at $x\in U_{g,\nu}$ is the space of \deleted{$g$} $u\in A_{g,\nu}^{\circ}$ such that $u(a_{g,\nu})=x$).
\inserted{$\unsure{=(A_{g,\nu}^{\circ})_{a_{g,\nu}}}$}
% Entre « de groupe » et A_{g,(\nu,\nu+1)}, une expression (A_{g,\nu})_{a_{g,\nu}} fortement surchargée est écrite dans la ligne ; une autre, =(A_{g,\nu})_{a_{g,\nu}}, en haut à droite au-dessus de « sur ».

% L13--16
The associated \inserted{Galois} covering, with group $B_{g,\nu}/B_{g,\nu}^{\circ}$, is therefore also the homogeneous space $A_{g,\nu}^{\circ}/A_{g,\nu+1}^{\circ}$, which is obviously connected, and pointed over $a_{g,\nu}\in U_{g,\nu}$. Put
\deleted{\uline{\unsure{Theorem}}}
\deleted{\uline{\illegible{2}}}
\margin{212}
\[
\tag{33}
\widetilde U_{g,\nu}=A_{g,\nu}^{\circ}/A_{g,\nu+1}^{\circ}.
\]

% L17--21, théorème et ajouts
\uline{Theorem}. $\widetilde U_{g,\nu}$ is simply connected \inserted{(if $(g,\nu)\ne(1,0)$ and $(g,\nu)\ne(0,2)$)} (and even contractible if $(g,\nu)\ne(0,0)$) and thus identifies with the universal covering pointed at $\unsure{a_{g,\nu+1}}$ of $U_{g,\nu}$.

% Marge gauche, à hauteur de L17--22
\margin{\unsure{This relation|This assertion} \unsure{plausible} \illegible{3} in the anabelian cases $2g+\nu\geq3$, \unsure{we shall see} \illegible{3} in abelian case \illegible{4} \unsure{reasonable} \illegible{3}. We must exclude the cases $(1,0)$ and $(0,2)$.}

% L22--23, formule (34), au bord inférieur
\uline{Corollary}. We have a canonical isomorphism
\[
\tag{34}
B_{g,\nu}/B_{g,\nu}^{\circ}\simeq\pi_1(U_{g,\nu};a_{g,\nu})\overset{\text{definition}}{=}\Pi_{g,\nu}.
\]

\scanpage{214}
\mspage{107}
% L1--4
Proof of the Theorem, in terms of ``well-known'' things. The \deleted{\illegible{1}} homotopy exact sequence for the fibration of $A_{g,\nu}^{\circ}$ over
\[
A_{g,\nu+1}^{\circ}/A_{g,\nu}^{\circ}\simeq\widetilde U_{g,\nu}
\]
\deleted{\unsure{gives}}
% L5, suite (35)
\[
\tag{35}
\longrightarrow\pi_{i+1}(\widetilde U_{g,\nu})
\longrightarrow\pi_i(A_{g,\nu+1}^{\circ})
\longrightarrow\pi_i(A_{g,\nu\deleted{\unsure{+1}}}^{\circ})
\longrightarrow\pi_i(\widetilde U_{g,\nu})
\longrightarrow\pi_{i-1}(A_{g,\nu+1}^{\circ})
\]
and it shows that $\pi_1(\widetilde U_{g,\nu})$ is isomorphic to the cokernel of
\[
\pi_1(A_{g,\nu}^{\circ})\longrightarrow\pi_1(A_{g,\nu+1}^{\circ}),
\]
whose kernel is a quotient of $\pi_2(\widetilde U_{g,\nu})$\ldots

% L9--13, formule (36)
If $(g,\nu)\ne0$, we know that $\pi_i(\widetilde U_{g,\nu})=0$ for $i\geq2$ (for every covering of $U_{g,\nu}$) hence
\deleted{\illegible{1}}
\[
\tag{36}
\left\{\begin{gathered}
\inserted{\text{if }(g,\nu)\ne(0,0)}\\
\pi_i(A_{g,\nu+1}^{\circ})\longrightarrow\pi_i(A_{g,\nu\deleted{\unsure{+1}}}^{\circ})
\text{ is }\begin{cases}
\text{bijective if }i\geq2\\
\text{injective with invariant image if }i=1
\end{cases}\\
\text{[the cokernel is }\simeq\pi_1(\widetilde U_{g,\nu})\text{]}
\end{gathered}\right.
\]
and we want to prove that it is even bijective for $i=1$. The thing to remember \unsure{easy} is this :

% L17--25, ensemble barré en diagonale
\deleted{\uline{Well-known theorem}.
$\pi_1(A_{g,\nu}^{\circ})=1$ \unsure{iff} We are in the anabelian cases (i.e. $2g+\nu\geq3$), \deleted{\illegible{1}}
\inserted{($\pi_1(U_{g,\nu})$ nonabelian i.e.\illegible{1})}
\deleted{Thus} [\inserted{This is equivalent} \deleted{\illegible{2}}, taking account of (35), \unsure{to} \deleted{we have}
\[
\begin{aligned}
\pi_1(A_{g,0}^{\circ})&=1\quad\text{if }g\geq2\\
\pi_1(A_{1,1}^{\circ})&=1\\
\pi_1(A_{0,3}^{\circ})&=1
\end{aligned}
\qquad
\begin{gathered}
\pi_1(A_{1,0}^{\circ}),\pi_1(A_{0,0}^{\circ}),\pi_1(A_{0,1}^{\circ})\\
\text{and }\pi_1(A_{0,2}^{\circ})\text{ are }\ne1.
\end{gathered}
\]
]}
\margin{213}

\scanpage{215}
% L1--8, théorème
\deleted{\unsure{and the} $\pi_1$}

\uline{Well-known theorem}. Equivalent conditions on the pair $(g,\nu)\in\mathbb N\times\mathbb N$:
\[
\begin{aligned}
\text{a)}\quad &2g+\nu\geq3\quad\text{i.e. we are not in one of the}\\
\text{\deleted{b)}}\quad &\text{four exceptional cases }(1,0),\ (0,0),\ (0,1),\ (0,2).\\
\text{b)}\quad &\pi_1(U_{g,\nu})\text{ nonabelian}\\
\text{c)}\quad &\deleted{\unsure{\pi}}\ A_{g,\nu}^{\circ}\text{ is simply connected}\\
&\inserted{\text{(}\deleted{\unsure{\text{and}}}\ (g,\nu)\ne(0,1)\text{)}}\\
\text{d)}\quad &A_{g,\nu}\text{ is contractible and }(g,\nu)\ne(0,1).
\end{aligned}
\]
% En d), A_{g,\nu} est écrit sans cercle (le cercle de c) est net).
\margin{abelian case}
% Cette note borde l'accolade à gauche de a) et b), à la hauteur des quatre couples exceptionnels.
\deleted{[\unsure{In the case} ``abelian'' \unsure{of} $U_{0,1}\simeq\mathbb C$ \unsure{simply} \illegible{3} $A_{g,\nu}^{\circ}$ \unsure{contractible}]}

% L9--14
In any case, whether or not these conditions are satisfied, we have \deleted{for $(g,\nu)$}
\[
\pi_i(A_{g,\nu}^{\circ})=0\qquad i\geq2,\quad(g,\nu)\ne0
\]
(i.e. \inserted{the} $A_{g,\nu}^{\circ}$ are classifying spaces of discrete groups $\unsure{\pi|K}$, with the sole exception of $A_{0,0}^{\circ}\simeq\operatorname{Aut}(S^2)^{\circ}$).

% L15--20, formules (37), (38)
Taking (36) into account, this Theorem is equivalent to the relations
\[
\tag{37}
\begin{aligned}
\pi_i(A_{g,0}^{\circ})&=0\quad\text{if }g\geq2,\ i\geq1
&&\text{i.e. for }g\geq2,\ A_{g,0}\text{ is contractible}\\
\pi_i(A_{1,1}^{\circ})&=0\quad\text{for }i\geq1
&&\text{i.e. }A_{1,1}^{\circ}\text{ is contractible}\\
\pi_i(A_{0,3}^{\circ})&=0\quad\text{for }i\geq1
&&\text{i.e. }A_{0,3}^{\circ}\text{ is contractible}.
\end{aligned}
\]
\[
\tag{38}
\left.\begin{aligned}
\pi_i(A_{1,0}^{\circ})&=0\quad\text{for }i\geq2\\
\pi_i(A_{0,1}^{\circ})&=0\quad\text{for }i\geq2
\end{aligned}\right\}
\quad\text{i.e. }A_{1,0}^{\circ}\text{ and }A_{0,1}^{\circ}\text{ are }K(\pi,1).
\]
\deleted{and it has \unsure{also} \unsure{as a consequence} that}

% L21--23, fin de page
and they \unsure{have} as a consequence that $\widetilde U_{g,\nu}$ is simply connected in the anabelian cases $2g+\nu\geq3$. To know how things stand
\margin{214}

\scanpage{216}
\mspage{108}
% L1--4
\unsure{in} the abelian cases, we must make precise the topological structure of $A_{g,\nu}^{\circ}$ in the \inserted{four} ``abelian'' cases $2g+\nu\leq2$. Now we have the

% L5--9, corollaire remanié
\uline{Corollary}. \deleted{\unsure{The conditions} of the Theorem} \inserted{For $U=U_{g,\nu}$ \deleted{\illegible{1}} \unsure{to have}} abelian $\pi_1$ (i.e. not satisfying the equivalent conditions of the theorem), \unsure{it is necessary and sufficient} that the topological space $U$ \deleted{be \unsure{isomorphic}} be able, \inserted{either} to be endowed with a topological group structure (which \unsure{will be} necessarily
\[
\simeq\mathbb U\times\mathbb U\quad\text{(case }(1,0)\text{)},\qquad
\simeq\mathbb C^*\quad\text{(case }(0,2)\text{)}\quad\text{or}\quad
\simeq\mathbb C\quad\text{(case }(0,1)\text{)}),
\]
\margin{$\unsure{\mathbb U}\overset{\text{definition}}{=}\{z\in\mathbb C^*\mid |z|=1\}$.}
% L10--13
or \inserted{with a structure} of a homogeneous space under a topological group \deleted{\illegible{1}}.
\inserted{(\unsure{we can take} $\operatorname{SO}(3,\mathbb R)$ \inserted{or $\unsure{\operatorname{GP}(1,\mathbb C)|\operatorname{GL}(1,\mathbb C)}$} for the case of $S^2$), by a connected subgroup.}

% L14--20
In all cases, the topological group \inserted{in question} can be described as follows: We choose a complex \deleted{\unsure{topology}} \inserted{structure} on $X_g=\widehat U_{g,\nu}=\widehat U$ (the \unsure{compactification} \unsure{pure} of $U$), and we take the induced complex structure \inserted{on $U$} (i.e. we choose an algebraic \inserted{complex} curve structure on $U$\ldots)
and we take $G\simeq\operatorname{Aut}_{\mathbb C}^{\circ}(U)$, \deleted{group} \unsure{the component} \deleted{connected} \inserted{of the identity} \unsure{of the group of} complex automorphisms of $U$, i.e. of the complex automorphisms of $\widehat U$ which \unsure{leave invariant|fix} $S_\nu=\widehat U\setminus U$.

% L21, fin de page
This being so, the inclusion $G\subset\operatorname{Aut}^{\circ}(U)\simeq A_{g,\nu}^{\circ}$
\margin{215}

\scanpage{217}
% L1--8, texte principal sous la note marginale
is a homotopy equivalence.

The fact that if $U$ is a homogeneous space of a group that is \inserted{topological, \unsure{by} a connected subgroup,} \deleted{connected}, then $\unsure{\pi_1(U)}$ is abelian, comes from the homotopy \unsure{exact} sequence, and from the fact that the $\unsure{\pi_1}$ of a topological group is \unsure{commutative}. \unsure{The} \unsure{converse} in the case of topological surfaces \illegible{1} remarkable! The corollary gives us

% Marge diagonale supérieure gauche
\margin{We have a more precise result: If $k$ is the smallest integer such that $(g,\nu+k)$ is an anabelian pair, then $\dim G=k$, and $G$ is simply transitive on the set of $k$-tuples $(x_1,\ldots,x_k)$ of distinct points of $U_{g,\nu}$, whence it follows at once that every element of $A_{g,\nu}$ can be written in a \unsure{unique} way as a product $gu$, with $g\in G$, $u\in\unsure{A_{g,(\nu,\nu+k)}}$, so \deleted{\illegible{2}} $\unsure{A_{g,\nu}}$ is homeomorphic to $A_{g,(\nu,\nu+k)}\times G$. As $G$ is connected, we conclude from this
\[
\Gamma_{g,\nu}\simeq\unsure{\Gamma_{g,(\nu,\nu+k)}}\subset\unsure{\Gamma_{g,\nu+k}},
\]
and $A_{g,\nu}^{\circ}\unsure{\simeq}G\times A_{g,\deleted{\illegible{1}}\unsure{\nu+}k}^{\circ}$, and as $A_{g,\nu+k}^{\circ}$ is \unsure{contractible|$\infty$-connected}, we conclude from this \illegible{2} $G\unsure{\xrightarrow{\approx}|\simeq}A_{g,\nu}^{\circ}$.}

% L9--14, tableau (39) et annotations à gauche
\[
\tag{39}
\begin{aligned}
\pi_1(A_{1,0}^{\circ})&\simeq\mathbb Z\times\mathbb Z
&&\text{i.e. }A_{1,0}^{\circ}\simeq K(\mathbb Z\times\mathbb Z,1)\\
\pi_1(A_{0,2}^{\circ})&\simeq\mathbb Z
&&\text{i.e. }A_{0,\unsure{\nu|2}}^{\circ}\simeq K(\mathbb Z,1)\\
\pi_1(A_{0,1}^{\circ})&\simeq\mathbb Z
&&\text{i.e. }A_{0,1}^{\circ}\simeq\deleted{\illegible{1}}\unsure{K(\mathbb Z,1)}.
\end{aligned}
\]
% Le numéro (39), entouré, est dans la marge gauche, sous la note diagonale.
\margin{$G\simeq\unsure{\mathbb U\times\mathbb U|U}$ (elliptic curve)}
\margin{$G\simeq\mathbb C^*=\operatorname{GL}(1,\mathbb C)$}
\margin{$G\simeq\operatorname{Aff}(1,\mathbb C)$ (homotopic to $\mathbb C^*$ by the inclusion)}
\margin{$G\simeq\operatorname{GP}(1,\mathbb C)$}

% L15--19
$A_{0,0}^{\circ}=\operatorname{Aut}(S^2)\simeq\operatorname{Aut}(\mathbb P_{\mathbb C}^1)$ is homotopic by the inclusion to its subgroup $\operatorname{GP}(1,\mathbb C)$ (hence to its maximal compact subgroup $\operatorname{SO}(3,\mathbb R)/\{\pm1\}$), which \uline{is not} a $K(\pi,1)$, and whose $\pi_1$ is $\simeq\mathbb Z/2\mathbb Z$.

% L20--23
We are going to use this corollary to determine the structure of $\widetilde U_{g,\nu}$ (as a Galois covering of $U_{g,\nu}$) for the ``abelian'' cases. Let us note for this that
\margin{216}

\scanpage{218}
\mspage{109}
% L1--5, calcul barré en diagonale
\deleted{1) Case $(g,\nu)=(1,0)$: elliptic surface $\simeq\unsure{\mathbb U\times\mathbb U}$ with $0$ holes, so $\pi_1(U_{g,0})\simeq\mathbb Z\times\mathbb Z$. In fact, I say that we even have in this case $B_{g,\nu}=B_{g,\nu}^{\circ}=A_{g,\nu+1}^{\circ}$ i.e. \deleted{$\unsure{B_{g,\nu}=}$} \uline{$B_{1,0}$ is connected}.}

% L6--8, formule (40)
the inclusion $G\subset A_{g,\nu}^{\circ}$ induces an inclusion $G_0\hookrightarrow B_{g,\nu}$, where \deleted{\illegible{1}}
\[
\tag{40}
G_0=G\cap B_{g,\nu}=\text{stabilizer of }a_{g,\nu}\text{ in }G.
\]

% L9--15
and as
\[
G/G_0\xrightarrow{\sim}A_{g,\nu}^{\circ}/B_{g,\nu}\simeq U_{g,\nu}
\]
($G$ being transitive on $U_{g,\nu}=U$) we find that the fact that $G\longrightarrow A_{g,\nu}^{\circ}$ is a homotopy equivalence, is equivalent to $G_0\longrightarrow B_{g,\nu}$ being one. Now in all the cases considered, $G_0$ is already connected: it is reduced to $1$ in the cases $(1,0)$ and $(0,2)$, [so in this case the corollary is equivalent to the \uline{contractibility} of $B_{g,\nu}$], and it is $\simeq\mathbb C^*$ in the case $(0,1)$, finally it is $\operatorname{Aff}(1,\mathbb C)$ in the case $(0,0)$. So we find
\margin{217}

\scanpage{219}
% L1--8, corollaire
\uline{Corollary}. In the abelian cases $2g+\nu<3$ (and only in those, of course) $B_{g,\nu}$ is connected (i.e. $=A_{g,\nu+1}^{\circ}$), i.e. $A_{g,\nu}^{\circ}/A_{g,\nu+1}^{\circ}\longrightarrow U_{g,\nu}$ is a homeomorphism. So $A_{g,\nu}^{\circ}/A_{g,\nu+1}^{\circ}\longrightarrow U_{g,\nu}$ makes the left-hand side a universal covering of $U_{g,\nu}$ iff $(g,\nu)$ is distinct from the two pairs $(1,0)$ and $(0,2)$.

% L9--14
Let us therefore assume $(g,\nu)\ne(1,0),(0,2)$ (morally we are working in the anabelian cases, the \deleted{\illegible{1}} included abelian cases $(0,1)$ and $(0,0)$ being \deleted{\illegible{1}} of no interest for what is to follow, it seems to me). Let us take up diagram (31) again, where in the second row the quotient group $B_{g,\nu}/B_{g,\nu}^{\circ}$ is therefore identified with:
\[
\Pi_{g,\nu}=\pi_1(U_{g,\nu},a_{g,\nu}).
\]
We therefore find

% L15--20, théorème, avec ajouts interlinéaires
\uline{Theorem}. \inserted{\unsure{Suppose} $(g,\nu)\ne(1,0)$ and $(0,2)$,} (the \unsure{subgroup} \inserted{$A_{g,(\nu,\nu+1)}/\unsure{A_{g,\nu}^{\circ}|A_{g,\nu+1}^{\circ}}$ of the \unsure{group} $=\Gamma_{g,\nu+1}'$} of Teichmüller
% Deux lignes interlinéaires au-dessus de « (le s-groupe … de Teichmüller » : la première porte la condition, la seconde le quotient.
\[
\Gamma_{g,\nu+1}=A_{g,\nu+1}/A_{g,\nu+1}^{\circ}
\]
\inserted{\uline{formed} of the \unsure{elements} that fix $a_{g,\nu}$}
admits a composition series \inserted{of length $3$}, whose successive factors are $\unsure{\mathfrak S_\nu|\mathfrak S_{\nu+1}}$ \deleted{(\illegible{2} $\unsure{\Gamma_{g,\nu+1}^{!}}$)}, $\Gamma_{g,\nu}^{!}$ and $\Pi_{g,\nu}$ --- \unsure{deduced} from a structure
\margin{218}

\scanpage{220}
\mspage{110}
% L1--3, formule (41)
of extension on $\Gamma_{g,\nu+1}^{!}$
\[
\tag{41}
1\longrightarrow\Pi_{g,\nu}\longrightarrow\Gamma_{g,\nu+1}^{!}
\longrightarrow\Gamma_{g,\nu}^{!}\longrightarrow1.
\]

% L4--8
[The outer operations of $\Gamma_{g,\nu}^{!}$ on $\Pi_{g,\nu}$ are those deduced by passage to the quotient from those of $A_{g,\nu}^{!}$ \inserted{$\subset\operatorname{Aut}(U_{g,\nu})$} \inserted{operating} by outer automorphisms on $\Pi_{g,\nu}=\pi_1(U_{g,\nu},a_{g,\nu})$.

% L9--15
Another way of saying things is this. Let us distinguish, in $\Gamma_{g,\nu+1}$ operating on $S_{\nu+1}$ \deleted{\illegible{1}} ($\simeq I_{\nu+1}$), the stabilizer of the last element $a_{g,\nu}$, i.e. of its complement, namely $\Gamma_{g,\nu+1}'$ \inserted{$\simeq A_{g,(\nu,\nu+1)}/A_{g,\nu+1}^{\circ}$}, \unsure{this|the} subgroup of index $\nu+1$, inverse image of $\unsure{\mathfrak S_\nu}$ \inserted{of $\unsure{\mathfrak S_{\nu+1}}$} under $\Gamma_{g,\nu+1}\longrightarrow\unsure{\mathfrak S_{\nu+1}}$.

% L16--19, formule (42)
This being so, we have an evident homomorphism of discrete groups, deduced from the inclusion \deleted{\illegible{1}} $A_{g,(\nu,\nu+1)}\hookrightarrow A_{g,\nu+1}$:
% Le but de l'inclusion porte l'indice \nu+1 dans l'encre (sic ; (42) demande A_{g,\nu}).
\[
\tag{42}
\Gamma_{g,\nu+1}'\longrightarrow\Gamma_{g,\nu}
\]
\margin{N.B. This homomorphism \unsure{still makes sense} without conditions on $(g,\nu)$ and it is an isomorphism in the abelian case $2g+\nu\leq2$.}
and this homomorphism is surjective (without conditions on $(g,\nu)$), finally (for $(g,\nu)\ne(1,0),(0,1)$) its kernel is canonically isomorphic to $\Pi_{g,\nu}$, so that we have an extension
\margin{\uline{N.B.}}

% Marge inférieure gauche, note barrée
\margin{N.B. \deleted{In the cases $(g,\nu)=\unsure{(0,2)}$, \illegible{2} abelian \illegible{3} $\Pi_{g,\nu}$ \illegible{1}, so $\Gamma_{g,\nu+1}'\unsure{\simeq}\Gamma_{g,\nu}$.}}

% L23, formule (43), au bord inférieur
\[
\tag{43}
1\longrightarrow\Pi_{g,\nu}\longrightarrow\Gamma_{g,\nu+1}'
\longrightarrow\Gamma_{g,\nu}\longrightarrow1,
\]
\margin{219}
